\chapter{Sirolimus Level}
\section*{What Is This Test?}
The sirolimus level (also known as rapamycin level) measures the whole blood concentration of sirolimus, an mTOR (mammalian target of rapamycin) inhibitor immunosuppressant used primarily in kidney transplantation as a calcineurin inhibitor-sparing agent. Sirolimus binds to FKBP12 (the same immunophilin as tacrolimus), but the sirolimus-FKBP12 complex inhibits mTOR Complex 1 (mTORC1), a serine/threonine kinase that regulates cell cycle progression from G1 to S phase. Inhibition of mTOR blocks IL-2-mediated signal transduction and T-cell proliferation. Sirolimus is also used in drug-eluting coronary stents to prevent restenosis, and in the treatment of lymphangioleiomyomatosis (LAM). Key advantages over calcineurin inhibitors: no nephrotoxicity. Key disadvantages: delayed wound healing (impaired fibroblast proliferation), hyperlipidemia/hypertriglyceridemia, myelosuppression (anemia, leukopenia, thrombocytopenia), proteinuria (glomerular podocyte injury), oral mucositis, interstitial pneumonitis, and impaired surgical wound healing (requiring withholding around transplantation surgery). Sirolimus has a very long half-life (57--63 hours). TDM uses whole blood (EDTA) collected as a trough (C0) level. Sirolimus is metabolized by CYP3A4/CYP3A5 and is a substrate of P-glycoprotein.
\section*{Alternative Names}
Rapamycin Level; Rapamune Level; Sirolimus Blood Level; Sirolimus Trough; mTOR Inhibitor Level
\section*{Principle}
Sirolimus is measured by LC-MS/MS (gold standard) or automated microparticle enzyme immunoassay (MEIA, Abbott IMx/ARCHITECT) in whole blood EDTA. Immunoassays cross-react with sirolimus metabolites. Specimen: whole blood EDTA (lavender-top). Trough: immediately before the next dose.
\section*{History}
Sirolimus (rapamycin) was isolated in 1975 from the bacterium Streptomyces hygroscopicus found in a soil sample from Rapa Nui (Easter Island), the source of its original name. It was first studied as an antifungal agent before its immunosuppressive properties were recognized in the late 1970s. In the 1990s, sirolimus was shown to inhibit the mammalian target of rapamycin (mTOR) pathway, and randomized trials demonstrated its efficacy in renal transplantation in combination with cyclosporine. The U.S. FDA approved sirolimus (Rapamune) for kidney transplant rejection prophylaxis in September 1999. Sirolimus-eluting coronary stents were introduced in the early 2000s, and the drug was approved for lymphangioleiomyomatosis (LAM) in 2015.
\section*{Why This Test Is Ordered}
Routine TDM in kidney transplant recipients. Target trough: 5--15 ng/mL (in combination with cyclosporine), 10--20 ng/mL (when cyclosporine is discontinued, sirolimus monotherapy or with mycophenolate). TDM essential because of significant inter-individual pharmacokinetic variability.
\section*{Primary vs Secondary Test}
This is a secondary, therapeutic-monitoring test rather than a diagnostic test. Sirolimus levels are monitored to keep whole blood trough concentrations within the target range, balancing the competing risks of acute rejection and drug toxicity.
\section*{How This Test Is Performed}
Blood is collected by standard venipuncture into an EDTA tube (lavender-top) because sirolimus distributes extensively into erythrocytes; serum or plasma is not used. A trough (C0) level is drawn immediately before the next scheduled dose, ideally at steady state (sirolimus half-life 57--63 hours, approximately 5--7 days after a dose change). The whole blood sirolimus concentration is measured by LC-MS/MS (gold standard) or by automated microparticle enzyme immunoassay (MEIA), with results reported in ng/mL. Results are typically available within 1--3 days for routine monitoring.
\section*{What Preparation Is Needed}
No fasting or special preparation is required. The level must be drawn as a true trough immediately before the next dose, and the dosing schedule should be consistent for at least 5--7 days before sampling to approach steady state. Record the time of the last dose, the dosing interval, and any recent changes in interacting drugs.
\section*{Contraindications and Precautions}
\begin{itemize}
  \item Recognize the signs of sirolimus toxicity: thrombocytopenia, anemia, hyperlipidemia, oral ulcers, proteinuria, interstitial pneumonitis, and impaired wound healing
  \item Sirolimus is a substrate of CYP3A4 and P-glycoprotein; strong CYP3A4 inhibitors (azole antifungals, macrolides) raise levels and inducers (rifampin, carbamazepine, phenytoin, St. John's wort) lower them --- recheck the level and adjust the dose after any interacting drug change
  \item Whole blood trough (C0) sampling is required; serum or plasma and random non-trough samples are not interpretable
  \item Immunoassays cross-react with sirolimus metabolites and may overestimate the concentration; use LC-MS/MS when results are inconsistent with the clinical picture
  \item Use caution when sirolimus is combined with cyclosporine, which increases sirolimus exposure; separate the doses or adjust per protocol
  \item Withhold or delay sirolimus around elective surgery because of impaired wound healing
\end{itemize}
\section*{Risks}
Minimal risk. Slight pain or bruising at the venipuncture site. Rarely: hematoma, infection, or vasovagal syncope.
\section*{What Happens After the Test}
The trough sirolimus level is interpreted against the target range (5--15 ng/mL with cyclosporine; 10--20 ng/mL without cyclosporine) in the context of the sampling time and clinical status. Subtherapeutic levels prompt a dose increase with rechecking after 5--7 days (to allow steady state), while supratherapeutic levels prompt dose reduction and evaluation for toxicity. Serum creatinine, fasting lipids, CBC, and urinary protein are followed in parallel to detect toxicity.
\section*{Understanding Results}
Subtherapeutic: Risk of acute rejection. Supratherapeutic: Myelosuppression (thrombocytopenia more common than anemia and leukopenia), hypertriglyceridemia and hypercholesterolemia, proteinuria, oral ulcers, delayed wound healing, interstitial pneumonitis, peripheral edema.
\section*{Normal Results}
Therapeutic trough (C0, whole blood EDTA): 5--15 ng/mL (with cyclosporine); 10--20 ng/mL (without cyclosporine). Half-life: 57--63 hours (steady state 5--7 days after dose change). TDM is essential at initiation, dose changes, drug-drug interactions, and periodically.
\section*{Follow-up Tests}
Fasting lipid profile, complete blood count, serum creatinine and urinary protein/spot urine protein-creatinine ratio, liver function tests, pulmonary function tests (if symptoms of pneumonitis), and drug interaction review.
\section*{References}
\begin{enumerate}
  \item Kahan BD. Sirolimus: a comprehensive review. Expert Opin Pharmacother. 2001;2(2):311--331.
  \item Augustine JJ, Bodziak KA, Hricik DE. Use of sirolimus in solid organ transplantation. Drugs. 2007;67(3):369--391.
  \item MacDonald AS. A worldwide, phase III, randomized, controlled, safety and efficacy study of a sirolimus/cyclosporine regimen. Transplantation. 2001;71(2):271--280.
  \item Kahan BD, Napoli KL, Kelly PA, et al. Therapeutic drug monitoring of sirolimus: correlations with efficacy and toxicity. Clin Transplant. 2000;14(2):97--109.
  \item Johnson RW, Kreis H, Oberbauer R, et al. Sirolimus allows early cyclosporine withdrawal in renal transplantation resulting in improved renal function and blood pressure. Transplantation. 2001;72(5):777--786.
\end{enumerate}
\clearpage
\section*{Regulatory Status and Authorization Date}
Regulatory status applies to the specific assay, instrument, manufacturer, and intended use rather than to the test name alone. No single approval date can be assigned to this test category. For a commercial assay, verify the current product in the relevant regulator's database: the U.S. Food and Drug Administration (FDA) clearance, approval, or authorization; India's Central Drugs Standard Control Organisation (CDSCO) licence or permission; the European Union In Vitro Diagnostic Regulation (EU IVDR) CE certificate issued through the applicable conformity-assessment route; or the United Kingdom Medicines and Healthcare products Regulatory Agency (MHRA) registration and UKCA/CE status. Laboratory-developed tests may not have an individual FDA, CDSCO, EU IVDR, or MHRA approval date.
\clearpage
